\documentclass[a4paper,10pt]{article}
\usepackage[margin=2cm]{geometry}
\usepackage{multicol}
\usepackage{times}
\usepackage{graphicx}

\title{ECG Heartbeat Classification Using Random Forest}
\author{Dinh Thanh Hai -- 22BA13119}
\date{}

\begin{document}
\maketitle

\begin{multicols}{2}

\section{Introduction}

Electrocardiogram (ECG) signals record the electrical activity of the heart and are widely used for diagnosing cardiovascular diseases. Different types of heartbeats correspond to different cardiac conditions, and accurate automatic classification of ECG beats is essential for clinical decision support systems.

Recently, deep learning models, especially Convolutional Neural Networks (CNNs), have achieved state-of-the-art performance in ECG analysis. These models can automatically learn complex features from raw signals, but they usually require large training datasets and high computational cost.

In this project, instead of using deep learning, a traditional machine learning approach based on Random Forest is applied to classify ECG heartbeats. The goal is to investigate whether a simpler model can still achieve high performance when trained on well-prepared ECG waveform data, and to compare the obtained results with a CNN-based method reported in the literature.

\section{Dataset Description}

The dataset used in this study is the MIT-BIH Arrhythmia Dataset (Kaggle version). The original MIT-BIH dataset contains long ECG recordings from multiple patients. In the Kaggle version, these recordings have been preprocessed and segmented into individual heartbeats.

Each heartbeat is represented by 187 normalized ECG samples, which correspond to one complete cardiac cycle. Therefore, each data sample is a vector of 187 real-valued features, followed by one class label.

The dataset contains five heartbeat classes:
\begin{itemize}
    \item Normal (N)
    \item Supraventricular Ectopic Beat (S)
    \item Ventricular Ectopic Beat (V)
    \item Fusion Beat (F)
    \item Unknown (Q)
\end{itemize}

The data is provided in two files: \texttt{mitbih\_train.csv} and \texttt{mitbih\_test.csv}. The training set is used to build the model, while the test set is used for performance evaluation. The dataset is highly imbalanced, with normal beats being the dominant class, while fusion and unknown beats appear much less frequently.

\section{Methodology}

A Random Forest classifier is used for ECG heartbeat classification. Random Forest is an ensemble learning method that builds a large number of decision trees during training. Each tree makes an independent prediction, and the final output is determined by majority voting.

The input to the model is the raw ECG waveform of each heartbeat, represented by 187 numerical values. No handcrafted features such as QRS duration or heart rate are used. This allows the model to directly learn discriminative patterns from the ECG shape.

Because the dataset is highly imbalanced, the parameter \texttt{class\_weight = balanced} is applied. This forces the classifier to give higher importance to minority classes during training, helping to improve the recognition of rare heartbeat types.

\section{Results}

After training, the Random Forest model was evaluated on the independent test set. The model achieved an overall classification accuracy of 97\%.

The performance on major heartbeat types was particularly strong:
\begin{itemize}
    \item Normal beats: F1-score = 0.98
    \item Ventricular Ectopic Beats: F1-score = 0.93
\end{itemize}

These results indicate that the model is able to correctly identify both healthy heartbeats and clinically important abnormal beats. Lower performance was observed for rare classes such as fusion beats, which is expected due to their small number of training samples.

\section{Comparison with the Original Paper}

The original research paper applies a deep convolutional neural network (CNN) to classify ECG heartbeats and evaluates the model according to the AAMI EC57 standard. Using this strict clinical protocol, the authors report an average accuracy of 93.4\% on the MIT-BIH Arrhythmia dataset.

In contrast, this project uses a Random Forest classifier and the Kaggle version of the dataset, where ECG beats are already segmented and provided in a fixed train--test split. Under this setting, the Random Forest achieved 97\% accuracy, which is higher than the accuracy reported by the CNN model.

However, this difference should be interpreted carefully. The AAMI EC57 standard uses patient-wise splitting and more realistic clinical evaluation, making the classification task more difficult. The Kaggle dataset, on the other hand, is easier because the signals are preprocessed and patient information is not considered.

Nevertheless, the high accuracy obtained in this project demonstrates that ECG waveforms in this dataset are highly separable and that even a traditional machine learning model can achieve very strong performance.

\section{Conclusion}

This project investigated ECG heartbeat classification using a Random Forest model on the MIT-BIH Arrhythmia dataset. Each heartbeat was represented by 187 ECG waveform values, allowing the classifier to directly learn from raw signal patterns.

Despite being a traditional machine learning approach without deep feature extraction, the Random Forest achieved 97\% accuracy on the test set, with excellent performance on both normal and abnormal heartbeats. This shows that well-prepared ECG data contains strong diagnostic information that can be captured even by relatively simple models.

Overall, this study confirms that while deep learning methods such as CNNs are powerful, traditional machine learning techniques like Random Forest remain highly effective for ECG classification tasks, especially when the data is clean and well-structured.

\end{multicols}
\end{document}
